\chapter{The Scoring Engine: Methodology and Implementation}
\label{chap:scoring-engine}

The two preceding chapters established \emph{what} this project set out to
build, a self-assessment platform for cyber crisis management maturity, grounded
in the ENISA \emph{Best Practices for Cyber Crisis Management} framework, and the
methodological principles that would govern \emph{how} it was built. This chapter
turns to the component on which everything else depends: the \emph{scoring engine},
the deterministic core that converts a respondent's answers into a maturity score.

The engineering problem is easy to state and deceptively hard to solve well. The
engine's input is a \emph{self-assessment}: for each of the framework's sixteen
best practices, a respondent declares a maturity level and points to the evidence
that supports it. Self-assessment is, by nature, subjective and declarative, two
properties that a scoring mechanism must actively counter rather than faithfully
reproduce. A naive approach, a weighted average of declared maturities, would fail
on three counts at once. It would be \textbf{opaque}: a single aggregate figure
gives a consultant no way to explain to a client \emph{why} the number is what it
is. It would be \textbf{gameable}: nothing would stop an organisation from claiming
the highest maturity everywhere, unsupported by any evidence, and walking away with
a flattering result. And it would be \textbf{fragile}: with nine best practices in
\emph{Preparedness} against only two in \emph{Response} and two in \emph{Recovery},
a plain mean would let one phase silently dominate the verdict.

Two properties are therefore non-negotiable, and they organise the entire design.
The score must be \textbf{auditable}: every point, every cap, every deduction
traceable back to a specific answer and a specific rule, so that the result can be
defended number by number in front of a client. And it must be \textbf{resistant to
gaming}: a claim unsupported by evidence must not be allowed to inflate the score.
From these two requirements the rest follows: rather than a single formula, the
engine is a pipeline of deliberately ordered stages, each one the answer to a
specific failure mode of naive scoring: evidence caps against the declarative
bias, foundational gates against structural weaknesses, consistency penalties
against incoherent self-portraits, and phase normalisation against structural
imbalance.

This chapter develops that design in five movements.
Section~\ref{sec:se-spec} derives the scoring specification from the consulting
need and fixes the engine's input and output contracts.
Section~\ref{sec:se-pipeline} presents the fixed-order pipeline at the heart of the
engine and argues why the order is itself a design decision.
Section~\ref{sec:se-architecture} shows how that specification is translated into a
codebase built for determinism and auditability.
Section~\ref{sec:se-testing} describes the testing strategy and an independent
audit that exposed, and led me to close, two genuine blind spots.
Section~\ref{sec:se-reflection} steps back to weigh the model's limits and the
choices I would defend or revisit.


\section{The scoring specification}
\label{sec:se-spec}

\subsection{From a consulting need to a specification}

The scoring specification has a concrete origin. The ENISA framework~\cite{enisa_crisis_mgmt_2024}
describes sixteen best practices across four crisis phases but provides no way to
measure adherence to them, the gap Chapter~\ref{chap:context} sets out in full. At the same time, the NIS2 directive~\cite{nis2_directive_2022} pushes organisations
to demonstrate crisis-management maturity, and the general-purpose self-assessments
that exist, such as the CISA Cyber Resilience Review~\cite{cisacrr} or maturity
ratings inspired by the NIST Cybersecurity Framework~\cite{nist_csf2_2024}, do not map
cleanly onto ENISA or onto NIS2 obligations. The need is therefore a measurement tool
that a Wavestone consultant can sit behind, in front of a client, to present a
maturity figure and defend it. That single constraint drives every requirement in
this section: a score that cannot be explained, or that rewards optimistic answers,
carries little value in an audit setting, whatever the formula behind it.

\subsection{The framework being scored}

The engine scores an organisation against the ENISA \emph{Best Practices for Cyber
Crisis Management} framework, encoded in a single file, \texttt{referential.yaml}
(version~1.1.0). This file is the centre of gravity of the whole specification: it
holds the four phases, the sixteen best practices, the forty-eight questions and
every threshold the engine uses. Table~\ref{tab:se-referential}, in Appendix~\ref{app:tables}, summarises its
structure. The framework follows the four phases of a crisis lifecycle, Prevention,
Preparedness, Response and Recovery, each carrying an equal weight of 25\% in the
global score. Every phase groups several best practices, and every best
practice is assessed through three questions.

For each question, a respondent provides two things: a declared maturity level on a
CMMI-style five-point scale~\cite{isaca_cmmi_cybermaturity}, from 0 (\emph{Not in place}) to 4
(\emph{Optimized}), and the
type of evidence backing the claim, from \texttt{no\_formal\_evidence} up to
\texttt{real\_incident}. The evidence type carries real weight, since it sets a
ceiling on how high the declared maturity can count. This is the first line of
defence against optimistic self-scoring. A best practice can also be marked
\emph{Not Applicable} when it genuinely does not fit the organisation, within strict
limits: at most three per assessment, never on a foundational best practice, and
always with a written justification of at least fifty characters.
Section~\ref{sec:se-pipeline} shows what that marking then does to the score.

\subsection{Rules carried by the referential}

Beyond the questions themselves, the referential also carries the methodological
adjustments that give the score meaning. Four best practices are flagged as
\emph{foundational} (BP04, BP07, BP08 and BP12), the governance and
crisis-coordination baseline the rest depends on. Three pairs of best practices are
declared as consistency pairs (BP10/BP11, BP07/BP08 and BP05/BP15), since the two
practices in a pair are expected to progress together. Each of these elements
carries its own thresholds: a minimum score and a cap for the gates, a maximum gap
and a point penalty for the pairs. Keeping these values in the referential rather
than in the code lets the framework evolve through a version bump, with no change
to the engine. Two further per-best-practice fields, \texttt{nis2\_\allowbreak articles} 
and \texttt{iso\_\allowbreak 27035\_\allowbreak clauses}, are declared alongside these rules 
but carry no values. They are empty placeholders, reserved for a future mapping of 
each best practice onto the directive's articles and the standard's clauses 
(Chapter~\ref{chap:impact}); no such mapping exists in version~1.1.0, and the engine 
never reads them.

\subsection{Output contract and requirements}

The output contract is where auditability is won or lost. The engine returns a full
result object, \texttt{AssessmentResult}, holding the final score together with the
reasoning behind it: which gates were applied, which consistency penalties were
deducted, which questions had their maturity capped by evidence, and which best
practices were excluded as N/A. A consultant can trace any figure back to the exact
rule and answer that produced it. The same object is defined once and reused by
every layer above the engine, and it later feeds the AI stage of
Chapter~\ref{chap:platform}, which reads the result and never modifies it.

These considerations translate into five requirements that the rest of the chapter
has to meet:
\begin{itemize}
  \item \emph{Auditability}: every point in the final score traces back to a rule
        and an answer.
  \item \emph{Anti-gaming}: unsupported or optimistic answers cannot inflate the
        score.
  \item \emph{Determinism}: the same input always produces the same output, with no
        hidden state.
  \item \emph{Reproducibility}: a score can be recomputed at any time from the
        stored answers and the referential version.
  \item \emph{Referential-driven}: the referential stays the single source of
        truth, and the engine reads its thresholds instead of hard-coding them.
\end{itemize}
Sections~\ref{sec:se-pipeline} and \ref{sec:se-architecture} show how the pipeline
and the code deliver each one.


\section{The fixed-order pipeline}
\label{sec:se-pipeline}

The engine computes a score in a fixed sequence of stages, shown in
Figure~\ref{fig:se-pipeline}. The order is itself part of the design, and this
section returns to it once the stages are in place. Each stage plays a precise
role: it either aggregates values up one level of the framework, from questions to
best practices to phases to the global score, or it applies one of the corrective
mechanisms that protect the two properties set out in Section~\ref{sec:se-spec}.
Read in order, the pipeline turns a set of raw answers into a single defensible
figure, and every intermediate value it produces is kept for the audit trail.

\subsection{Capping declared maturity by evidence}

The first stage acts on each question, before any aggregation. The declared
maturity is capped by the evidence type: a level claimed with
\texttt{no\_formal\_evidence} cannot count above 1, \texttt{documented} above 2,
\texttt{tested} above 3, and only \texttt{real\_incident} allows the full 4. The
effective maturity used from then on is the smaller of the declared level and this
cap. This stage is the engine's answer to the declarative bias of self-assessment:
a confident claim with nothing to back it up is brought back to what the evidence
supports, and the cap is recorded so the respondent can see why. The question score
is then the effective maturity over 4, weighted by the question's own weight.

\subsection{Aggregating to best practices and phases}

The next stages aggregate upward. Question scores combine into a best-practice score
on a 0 to 100 scale, and best-practice scores combine into a phase score using their
intra-phase weights. A best practice marked N/A leaves the computation entirely,
with no contribution to either the numerator or the denominator of its phase, so a
genuinely irrelevant practice neither helps nor hurts the result. Phase scores then
combine into the raw global score, where each of the four phases counts for exactly
25\%, whatever the number of best practices it contains. This last choice answers
the structural imbalance noted earlier: left alone, Preparedness and its nine best
practices would speak far louder than Response or Recovery with two each.

\subsection{Foundational gates}

With a raw global score in hand, the engine applies the foundational gates. If a
foundational best practice (BP04, BP07, BP08 or BP12) scores below its required
minimum of 40, the global score is capped at the value set for that gate, between 60
and 70 depending on the practice. When several gates apply at once, the strictest
cap wins. The gates encode a simple idea: some practices are prerequisites, and a
serious weakness in governance or coordination should hold the whole score down,
however strong the rest looks. A high average cannot paper over a missing
foundation.

\subsection{Consistency penalties and finalisation}

The engine then applies the consistency penalties, the most original mechanism of
the model. Three pairs of best practices are expected to move together: exercises
and training (BP10/BP11), roles and escalation criteria (BP07/BP08), and asset
mapping and business resumption (BP05/BP15). For each pair, the engine measures the
gap between the two effective maturities, averaged over their questions. A pair is
skipped whenever either of its two best practices is marked N/A, since a practice
that does not apply cannot be called inconsistent with its partner. If the gap
exceeds two levels, it deducts a flat five points from the score. The penalty
targets incoherent profiles, where an organisation rates one practice highly while
neglecting the practice that should support it. Penalties are cumulative, and the
score is floored at 0. A final stage maps the resulting score onto the five global
maturity levels, from \emph{Not in place} to \emph{Optimized}, for reporting to an
executive audience. Table~\ref{tab:se-levels}, in Appendix~\ref{app:tables}, gives the bands, which tile the whole
0 to 100 range with no gap and no overlap.




\subsection{A worked example}
The example below follows a single assessment through the pipeline. It is
deliberately simple, chosen so that one run triggers an evidence cap, a foundational
gate and a consistency penalty, with each mechanism visible in turn.
Table~\ref{tab:se-worked-example} gives the running value after each stage.

\begin{table}[htbp]
  \centering
  \caption{End-to-end worked example: from raw answers to the final global score.}
  \label{tab:se-worked-example}
  \begin{tabular}{@{}p{2.8cm}p{3.0cm}p{5.4cm}@{}}
    \toprule
    Pipeline stage & Running value & Why \\
    \midrule
    Evidence cap           & BP08\_Q3: 3 capped to 1 & No formal evidence limits the claim to \emph{Ad hoc}. \\
    BP / phase aggregation & BP08 $= 30/100$         & Weak effective maturities across its three questions. \\
    Phase normalisation    & Raw global $= 72$       & Mean of the four phase scores, 25\,\% each. \\
    Foundational gate      & Capped to $65$ (BP08 $< 40$) & A weak escalation baseline caps the whole score. \\
    Consistency penalty    & $65 - 5 = 60$           & BP07 and BP08 differ by more than two maturity levels. \\
    Finalisation           & $60 \rightarrow$ Level~2 & Mapped to \emph{Defined} for executive reporting. \\
    \bottomrule
  \end{tabular}
\end{table}


\subsection{Why gates come before penalties}

The order of the last two mechanisms is a deliberate choice, and it changes the
result. In the example just worked through, applying the gate first caps the raw 72
at 65, then the penalty brings it to 60. Reversing the order would
deduct the penalty first, giving 67, then cap at 65, so the incoherence would
disappear into the ceiling. The chosen order treats the gate as a hard ceiling caused
by a missing prerequisite and makes the incoherence penalty bite on top of it, which
is the stricter and more faithful reading. The same logic explains why the strictest
cap wins when two gates apply: the ceiling reflects the worst prerequisite failure,
not an average of them.

\begin{figure}[htbp]
  \centering
  \begin{tikzpicture}[
      >={Latex[length=2mm]},
      node distance=6.5mm,
      stage/.style={draw, rounded corners, align=center, text width=4.4cm,
                    minimum height=8.5mm, font=\small},
      guard/.style={stage, fill=gray!15, very thick},
      note/.style={align=left, text width=4.6cm, font=\footnotesize\itshape,
                   text=gray!50!black},
      arr/.style={-{Latex[length=2mm]}, thick},
    ]
    \node[stage] (input)             {Answers: declared maturity $+$ evidence};
    \node[guard, below=of input] (cap)  {Evidence cap \\ (per question)};
    \node[stage, below=of cap]   (q)    {Question score};
    \node[stage, below=of q]     (bp)   {Best-practice score};
    \node[stage, below=of bp]    (ph)   {Phase score};
    \node[stage, below=of ph]    (gl)   {Raw global score \\ (4 phases, 25\,\% each)};
    \node[guard, below=of gl]    (gate) {Foundational gates};
    \node[guard, below=of gate]  (pen)  {Consistency penalties};
    \node[stage, below=of pen]   (fin)  {Final score \& maturity level};

    \draw[arr] (input) -- (cap);
    \draw[arr] (cap)   -- (q);
    \draw[arr] (q)     -- (bp);
    \draw[arr] (bp)    -- (ph);
    \draw[arr] (ph)    -- (gl);
    \draw[arr] (gl)    -- (gate);
    \draw[arr] (gate)  -- (pen);
    \draw[arr] (pen)   -- (fin);

    \node[note, right=7mm of cap]  {counters the declarative bias};
    \node[note, right=7mm of gl]   {counters structural imbalance};
    \node[note, right=7mm of gate] {enforces prerequisites; strictest cap wins};
    \node[note, right=7mm of pen]  {penalises incoherent pairs};
  \end{tikzpicture}
  \caption{The fixed-order scoring pipeline. Each stage answers a specific
           failure mode of a naive weighted average.}
  \label{fig:se-pipeline}
\end{figure}


\section{Architecture of the engine}
\label{sec:se-architecture}

\subsection{Two principles, expressed in code}

The specification of Section~\ref{sec:se-spec} and the pipeline of
Section~\ref{sec:se-pipeline} translate into a codebase whose shape is driven by two
of the inviolable principles set out in Chapter~\ref{chap:method}: determinism (P1)
and the referential as single source of truth (P4). The engine is a pure,
synchronous function with no I/O and no randomness. The same input always produces
the same output. That property is what turns reproducibility and auditability from
intentions into facts: any score can be recomputed at will, and there is no hidden
state to explain away when defending a figure.

\subsection{One module per pipeline stage}

The code mirrors the pipeline one module at a time. A \texttt{referential\_loader}
reads and validates the YAML, an \texttt{evidence} module applies the caps, an
\texttt{na\_rules} module handles the N/A cases, an \texttt{aggregation} module
climbs from questions to best practices to phases, and a \texttt{finalization}
module applies the gates, the penalties and the level mapping. A thin
\texttt{engine} module orchestrates them in the fixed order. Each function is pure
and short, so the top-level routine reads like the ordered list of
Section~\ref{sec:se-pipeline} rather than hiding it under machinery.
Listing~\ref{lst:se-engine} in Appendix~\ref{app:figures} shows that routine, and
Figure~\ref{fig:se-modules}, in the same appendix, shows the decomposition.

\subsection{Three families of data models}

The data models are split into three families that the engine keeps deliberately
apart: the referential models, which hold the framework loaded from the YAML; the
input models, which hold the raw answers a respondent provides; and the output
models, built around the \texttt{AssessmentResult} audit trace. The separation keeps
three questions from ever being confused: what the framework says, what the user
claimed, and what the engine concluded. It also anticipates the platform, where the
database stores the raw inputs while the result stays recomputable on demand.

\subsection{Defensive design and guaranteed bounds}

Defensive design runs through the whole package. The referential models are frozen
and reject unknown fields, so the framework cannot be mutated after loading and a
typo in the YAML raises an error instead of quietly creating a phantom field. The
loader checks its invariants at startup: every best practice carries exactly three
questions, the global maturity levels tile the whole 0 to 100 range with no gap, and
no foundational practice is allowed to be N/A. A malformed referential fails loudly
and early, before a single assessment is scored.

The final score is bounded to the range 0 to 100 twice over, once by the arithmetic
of the pipeline and once by a constraint on the output field, so a computation bug
surfaces as a validation error rather than as an implausible score reaching a
client. The package is fully typed throughout. Taken one by one these guarantees are
modest, but together they make the engine hard to break quietly, which is exactly
what a component whose output has to be trusted in front of a client needs.

\subsection{A deliberately isolated core}

The engine stays deliberately unaware of everything around it: no database, no web
framework, no configuration beyond the path to the referential. When the platform is
built (Chapter~\ref{chap:platform}), a single translation layer maps stored answers
onto the engine's input models, and the score is computed on demand from those
answers and the referential version. Keeping that seam thin and one-directional is
what lets the engine remain the pure, testable core that the next section puts under
scrutiny.


\section{Testing strategy and audit}
\label{sec:se-testing}

\subsection{The test suite}

An engine is only as trustworthy as the tests that pin its behaviour down, and here
the tests carry an unusual weight, since the whole point of the engine is to be
defended figure by figure. The suite is organised by pipeline stage, so each
mechanism is checked in isolation, and it is anchored by five \emph{golden}
scenarios that act as an executable specification: a fixed set of answers mapped to
a fixed, fully written-out \texttt{AssessmentResult}. A dedicated determinism test
enforces P1 by scoring the same input twice and comparing the serialised results
byte for byte. A parametrised sweep over many answer combinations proves that the
score always lands within 0 to 100, and the loader's failure cases (a missing file,
malformed YAML, a missing field) are covered as well. At the point of the audit the
suite held 64 tests, all green.

\subsection{An independent audit and its blind spots}

To get an outside view before any platform work began, the engine, the referential
and the documentation were put through an independent, read-only audit. It returned
five findings, and the two rated high severity are what made it worth doing, two
real blind spots that a set of green tests had kept out of sight.

The first and sharpest finding (H1) concerned the consistency penalty, the model's
most original mechanism. No test ever exercised it. Every golden scenario used a
uniform maturity across best practices, so the maturity gap never crossed the
threshold and the penalty branch simply never ran. A suite of green tests told a
reassuring story while leaving the one mechanism that makes the model distinctive
entirely uncovered. The multi-gate rule, where the strictest cap wins, was checked
only loosely, without a numerical assertion.

The second finding (H2) was a naming contradiction. The penalty field was called
\texttt{penalty\_percent\_\allowbreak of\_\allowbreak global\_score}, yet the code
subtracted it as a flat number of points. On a score of 80, five percent and five points are not the same
value, and a later maintainer could easily have corrected the code in the wrong
direction. The audit also flagged two referential fields that were loaded and
validated but never actually read (\texttt{global\_weight} and
\texttt{foundational\_can\_be\_na}), along with an undocumented rounding convention.
Table~\ref{tab:se-audit} lists the findings and how each was closed.

\subsection{The corrective loop}

The corrections followed a deliberate order, and the order is the point. A safety
net came first: three targeted tests were written before any production code was
touched, capturing the current behaviour of the penalty branch, the strictest-cap
rule, and the strict greater-than threshold. Only once that behaviour was
locked did the refactors follow. The penalty field was renamed to
\texttt{penalty\_points} and its meaning fixed as a flat, absolute deduction, which
changes the YAML contract and so triggered a referential version bump from 1.0.0 to
1.1.0. The two decorative fields were wired to drive the computation for real, each
change designed to leave the golden numbers untouched when the phase weights are all
equal, so the existing scenarios stayed green. The rounding convention was
documented rather than altered. The suite was re-run after every step and moved from
64 to 67 tests, all passing, with no existing test modified.

This sequence is what the chapter has been building toward. Writing the safety net
before the refactor, bumping the version when the contract changed, and preferring
changes that preserve the golden numbers are the habits that let a scoring engine be
modified without fear of a silent regression. The audit turned ``the tests
pass'' into ``the differentiator is proven'', which is the claim that actually
matters in front of a jury or a client. After integration into the platform the
suite stands at 70 tests. What it cannot settle, the choices that are defensible
rather than provable, is the subject of the next section.

\begin{table}[htbp]
  \centering
  \caption{Audit findings and the corrective loop that followed.}
  \label{tab:se-audit}
  \begin{tabular}{@{}cp{4.7cm}cp{5.6cm}@{}}
    \toprule
    ID & Finding & Severity & Resolution \\
    \midrule
    H1 & Consistency-penalty branch never exercised by tests & High &
         Three golden tests added; suite $64 \rightarrow 67$ \\
    H2 & Field name contradicts the absolute-points code & High &
         Renamed \texttt{penalty\_points}; referential bumped to 1.1.0 \\
    M1 & \texttt{global\_weight} loaded but unused & Medium &
         Weighted mean now reads it (invariant at 0.25) \\
    M2 & \texttt{foundational\_can\_be\_na} inert & Medium &
         Flag now read from the referential \\
    B1 & Rounding convention undocumented & Low &
         Documented (round-half-to-even) \\
    \bottomrule
  \end{tabular}
\end{table}


\section{Critical reflection}
\label{sec:se-reflection}

\subsection{Defensible, not provable}

The engine is deterministic and has been audited, but a few of its choices are
matters of expert judgement, defensible rather than provable. Setting them out
plainly is part of the same engineering honesty that drove the audit, and it is also
the best preparation for the questions a jury is likely to ask.

The clearest example is the set of thresholds the model relies on: the minimum score
of 40 that a foundational practice must reach, the caps of 60 to 70 that the gates
apply, the maximum maturity gap of two, and the five-point penalty. These values
come from Wavestone expert judgement rather than from a calibration on data. They
live in the referential, so any of them can be adjusted through a version bump, but
their exact figures are the weakest empirical point of the model. A natural next step
would be a sensitivity analysis, showing how much the global score moves when a cap
or a threshold is shifted, so that the numbers can be discussed on evidence instead
of on authority alone.

\subsection{Measuring demonstrable maturity}

A second choice worth stating openly is that the consistency penalty compares
effective maturity, measured after the evidence cap. A practice that is rated highly
but backed by weak evidence is capped, and can therefore trigger a penalty against a
better-evidenced neighbour. This is deliberate: the model rewards maturity that can
be demonstrated rather than merely declared. It deserves to be said out loud, because
comparing declared maturity instead is equally easy to implement and would simply
need a different justification.

\subsection{Smaller deliberate choices}

A few smaller decisions are of the same kind. The four phases carry equal weight, the
structural choice argued in Section~\ref{sec:se-pipeline}; that weighting is now read
from the referential, so a later version could revisit it. The penalty threshold is strict, so a gap of exactly
two levels does not penalise while a gap just above it does. The global score is
rounded to the nearest even value at the level boundaries. None of these is
controversial on its own, yet each is a decision a careful reader could question, so
each is documented and covered by a test instead of left implicit.

\subsection{A caveat on N/A}

One behaviour deserves a direct caveat for anyone reading a generated report. Since
an N/A best practice leaves its phase entirely, marking both practices of a short
phase as N/A makes that phase drop out of the global score rather than scoring zero.
That is consistent with treating N/A as neutral, but it means a result has to be read
as a phase that was not assessed, rather than a phase that scored nothing, and the
interface has to carry that distinction so a reader is never misled.

\subsection{What the engine enables next}

Set against these limits, the two properties the chapter opened with are what give
the engine value beyond this single project. Because the score is deterministic and
fully traced, the AI layer added in Chapter~\ref{chap:platform} can be held to a
strict and safe role: it reads the \texttt{AssessmentResult} and the referential to
phrase recommendations for the weakest best practices, never client names or the
private N/A justifications, and any failure on its side degrades to a deterministic
recommendation drawn from the referential (Chapter~\ref{chap:platform}). The engine remains the single arbiter of
the numbers, and the parts of the model that stay matters of judgement are exactly
the ones Chapter~\ref{chap:pm} returns to with the benefit of hindsight.


% =====================================================================
%  Fin du chapitre 5
% =====================================================================